% !TeX program = xelatex
% Encoding: UTF-8
% Compile: xelatex -no-shell-escape -interaction=nonstopmode -halt-on-error liutao-resume.tex
% Requires XeLaTeX, fontspec, geometry and hyperref.
% Chinese fonts: Heiti SC (macOS), Noto Sans CJK SC, or Fandol (TeX Live).
\documentclass[a4paper,11pt]{article}

\usepackage[margin=17mm]{geometry}
\usepackage{fontspec}
\usepackage[hidelinks]{hyperref}

% Select an available font without requiring ctex or xeCJK.
\IfFontExistsTF{Heiti SC}{
  \setmainfont{Heiti SC}[AutoFakeBold=2]
}{
  \IfFontExistsTF{Noto Sans CJK SC}{
    \setmainfont{Noto Sans CJK SC}
  }{
    \setmainfont{FandolHei-Regular.otf}[BoldFont=FandolHei-Bold.otf]
  }
}
\XeTeXlinebreaklocale "zh"
\XeTeXlinebreakskip=0pt plus 1pt
\pagestyle{empty}
\setlength{\parindent}{0pt}
\setlength{\parskip}{4pt}
\linespread{1.08}
\renewenvironment{itemize}{
  \begin{list}{\textbullet}{
    \setlength{\leftmargin}{1.1em}
    \setlength{\labelwidth}{0.7em}
    \setlength{\labelsep}{0.4em}
    \setlength{\itemsep}{4pt}
    \setlength{\topsep}{4pt}
    \setlength{\parsep}{0pt}
    \setlength{\partopsep}{0pt}
  }
}{\end{list}}
\makeatletter
\renewcommand\section{\@startsection{section}{1}{0pt}{-10pt}{5pt}{\large\bfseries}}
\renewcommand\subsection{\@startsection{subsection}{2}{0pt}{-7pt}{3pt}{\normalsize\bfseries}}
\makeatother
\setcounter{secnumdepth}{0}
\emergencystretch=1.5em
\widowpenalty=10000
\clubpenalty=10000
\hypersetup{
  pdftitle={刘韬 前端架构与研发效能简历},
  pdfauthor={刘韬},
  pdfsubject={职业经历与代表项目},
  pdfkeywords={前端,跨端SDK,研发效能,互动白板,Lynx}
}

\newcommand{\role}[2]{\textbf{#1}\hfill #2\par}
\newcommand{\lead}[1]{\textbf{#1：}}

\begin{document}

{\fontsize{26pt}{31pt}\selectfont\bfseries 刘韬}\par
前端架构与研发效能方向\quad |\quad
\href{mailto:liutao.fe@bytedance.com}{liutao.fe@bytedance.com}

2020 年 10 月加入字节跳动，历经商业化、视频云与生活服务。具备复杂前端产品、跨端 SDK 和研发工具建设经验，曾担任互动白板项目 Owner，推动产品能力、性能、工程质量与客户接入的持续改进。

\section{工作经历}

\role{字节跳动 · 生活服务交易前端}{2024.11 至今}
交易业务与研发效能
\begin{itemize}
  \item \lead{交易业务与架构}参与订单、提单等交易链路研发，支持视频跟买、通兑链路优化、旅行社出团通知书线上化、商家会员和学生专属商品等需求；完成订单预约卡二期标准化改造并进入实验。
  \item \lead{资产采集与调试}建设 DitoM 交易资产插件相关能力，串联真机画面、接口抓取、Mock 数据、截图和测试商品；集成任意门代理能力，支持场景留存与复现，并上线测试商品查找工具。
  \item \lead{资产质量治理}负责卡片截图采集、运行时标识关联与采集链路优化，支持订单、提单及商详场景。2026 Q2 团队复盘记录资产准确率由 \textbf{81.81\% 提升至 92.31\%}、场景与 UI 覆盖率由 \textbf{38.71\% 提升至 58\%}，平台用于内部试用。
  \item \lead{Oncall 流程改进}推进服务端、客户端和前端值班租户融合与标签对齐，沉淀排查 SOP，并引入 Aime 辅助排查；2025 Q2 个人复盘记录团队平均工单时长由 \textbf{404.01 小时降至 131 小时}。
\end{itemize}

\role{字节跳动 · 视频云前端}{2022.10 至 2024.10}
互动白板与 WebRTC 解决方案
\subsection{跨端 SDK 自动生成}
\begin{itemize}
  \item 针对跨端 SDK 大量重复封装的问题，建设从 Native API 信息到 TypeScript SDK 的生成链路：基于 Doxygen 二次开发提取 API，转换为结构化数据并生成代码，通过 IPC 与 Native 反射完成运行时调用。
  \item 2024 Q2 在直播 React Native SDK 落地：拉流覆盖 \textbf{69 个 API}，投入由人工方案预估 \textbf{45 人天降至 13 人天}；推流覆盖 \textbf{242 个 API}，由预估 \textbf{50 人天降至 19 人天}。
\end{itemize}

\newpage

\section{视频云代表项目}
\subsection{火山互动白板}
项目 Owner\quad |\quad 主要建设期 2022.10 至 2024.03，后续持续维护
\begin{itemize}
  \item \lead{产品与架构}负责整体技术设计和产品能力建设，覆盖实时绘制、多人协作、录制回放、动态 PPT 与多端接入。采用 Canvas 2D、实时信令和 Protobuf，结合元素锁、数据一致性检测及快照机制支撑协同与回放。
  \item \lead{跨端复用}主导移动端 SDK 从 Native 架构迁移到 WebView 与 Hybrid 架构，实现一套核心代码三端共享；通过通用实例与方法调用协议减少桥接层适配，单需求 Native 人力投入由 \textbf{3 降至 0.5}。
  \item \lead{性能与可靠性}推动消息传输链路优化，端到端延迟由 \textbf{600ms 降至 100ms}；参与转码架构和重试机制优化，静态转码最终成功率由 \textbf{96\% 提升至 99\%}，动态转码成功率由 \textbf{82\% 提升至 98\%}。
  \item \lead{工程质量}引入并持续完善单元测试和 CI，代码覆盖率由 \textbf{0 提升至 74\%}，千行 Bug 率由 \textbf{30+ 降至 1 以下}；建设 20 余个数据看板，覆盖消息、转码和关键交互指标。
  \item \lead{客户落地}支持企业直播接入，通过云端无头浏览器采集白板与动态 PPT 并合流推送；优化 SDK API，使初始化代码由约 \textbf{100 行减少至 10 行以内}。参与 WebRTC 售前体验、接入实践与 Oncall，支持内部 10 余个、外部 9 余个客户。
\end{itemize}

\section{商业化与早期经历}
\role{字节跳动 · 商业化落地页前端}{2020.10 至 2022.10}
橙子建站与广告自助创意
\begin{itemize}
  \item 完成移动端落地页编辑器的设计与开发，并在 DOU+ 落地；衔接模板管理、组件资源、Node.js BFF 与落地页渲染，支持广告主基于模板编辑和投放。项目记录日均有消耗站点 \textbf{87 个}、日均广告消耗 \textbf{1.7 万+}，管理模板 \textbf{300 余个}。
\end{itemize}

\role{携程 · Trip.com 国际事业部公共前端}{2018.06 至 2020.10}
负责 Foxpage 页面搭建平台的 C 端渲染与组件开发工具。平台支持 95\% 的国际业务营销页面搭建，一年上线 1,200 余个页面，营销页面开发时间由 40 小时缩短至 2 小时。

\section{教育背景}
\textbf{四川大学}\quad 本科 · 信息管理与信息系统\hfill 2017 年毕业

\section{技术能力}
TypeScript / JavaScript、Node.js、Lynx、React Native、WebView / JSBridge；Canvas 2D、WebRTC / RTS、Protobuf；SDK 设计与代码生成、页面搭建、单元测试与 CI、真机调试、Mock 与资产采集。

\end{document}
